%!TEX TS-program = pdflatex
\documentclass[11pt,a4paper]{article}
\usepackage[T1]{fontenc}
\usepackage[utf8]{inputenc}
\usepackage[ngerman]{babel}
\usepackage{amsmath,amssymb}
\usepackage[margin=22mm]{geometry}
\usepackage{enumitem}
\usepackage{fancyhdr}
\newif\ifloesungen
\loesungentrue % Für die Ausgabe ohne Musterlösungen: \loesungenfalse
\pagestyle{fancy}
\fancyhf{}
\fancyhead[L]{Theoretische Informatik}
\fancyhead[R]{Formale Sprachen und Automaten}
\fancyfoot[C]{\thepage}
\setlength{\headheight}{14pt}
\setlength{\parindent}{0pt}
\setlength{\parskip}{5pt}
\setlist[enumerate]{leftmargin=*,itemsep=6pt}
\renewcommand{\arraystretch}{1.15}
\newcommand{\punkte}[1]{\hfill\mbox{\textbf{(#1 Punkte)}}}
\begin{document}
\begin{center}
{\LARGE\bfseries Klausur: Formale Sprachen und Automaten}\\[6pt]
{\large\ifloesungen Aufgaben und separate Musterlösungen\else Aufgaben\fi}\\[5pt]
Bearbeitungszeit: 90 Minuten \quad Gesamtpunktzahl: 60
\end{center}
\textbf{Name:} \hrulefill\quad\textbf{Matrikelnummer:} \hrulefill

\textbf{Hinweise.} Bearbeiten Sie alle drei Aufgaben auf separatem Papier.
Begründen Sie Ihre Antworten und dokumentieren Sie die Rechenschritte.
DFA sind vollständig anzugeben. Bei NFA dürfen Übergangsmengen leer sein.
Ein Tupel mit Übergangstabelle genügt; ein Zustandsgraph ist nicht erforderlich.
$\to$ kennzeichnet den Startzustand, $*$ einen Endzustand.
$\mathbb N_0=\{0,1,2,\ldots\}$.

\section*{Aufgabe 1: Von der Sprache zum minimalen DFA \hfill 30 Punkte}
Gegeben sind $\Sigma=\{a,b\}$ und die Sprache
\[
 L=\{w\in\Sigma^*\mid w\text{ endet auf }aba,\ bba\text{ oder }bab\}.
\]
Gemeint sind die \emph{letzten drei Zeichen}: Ein bloßes Vorkommen eines
Musters im Wort reicht nicht aus. Wörter der Länge kleiner als drei gehören
nicht zu $L$.

\begin{enumerate}[label=\alph*)]
\item \textbf{NFA konstruieren.} Bauen Sie einen NFA ohne
$\varepsilon$-Übergänge für $L$ nach folgendem Konstruktionsprinzip:
Verwenden Sie einen Startzustand $q_0$, in dem ein beliebiges Präfix gelesen
werden kann. Für jedes der drei Endmuster soll vom Startzustand eine
\emph{eigene Kette mit drei Übergängen} abzweigen. Verwenden Sie für diese
Ketten jeweils drei neue Zustände; teilen Sie keine Kettenzustände und
fassen Sie noch keine Zustände zusammen.

Geben Sie das vollständige Tupel und die Übergangstabelle an.
Erklären Sie den Nichtdeterminismus sowie kurz, warum Ihr NFA genau $L$
erkennt. \punkte{10}

\item \textbf{DFA konstruieren.} Wandeln Sie Ihren NFA aus a) mit der
Potenzmengenkonstruktion in einen vollständigen DFA um. Geben Sie alle
\emph{erreichbaren} Zustandsmengen, ihre Übergänge, Start- und Endzustände an.
Nennen Sie für jede Zustandsmenge ein Wort, durch das sie erreicht wird.
Rechnen Sie einen Übergang als Vereinigung von Übergangsmengen vor.
Minimieren Sie hier noch nicht. \punkte{12}

\item \textbf{DFA minimieren.} Minimieren Sie den DFA aus b) mit dem
Tabellenalgorithmus. Dokumentieren Sie die Anfangsmarkierungen und die
weiteren Markierungsrunden in der Zustandspaartabelle. Geben Sie die
Äquivalenzklassen sowie die Übergangstabelle des minimalen DFA an.
Begründen Sie Erreichbarkeit und Minimalität. \punkte{8}
\end{enumerate}

\textit{Hinweis:} Die Bauvorgabe in a) sorgt für eine vergleichbare Ausgangsbasis
für b) und c). Nichtdeterministische Entscheidungen dürfen auch bei
gleichem Eingabezeichen vom selben Zustand ausgehen.

\newpage
\section*{Aufgabe 2: Pumping-Lemma \hfill 12 Punkte}
Betrachten Sie die Sprache über dem einelementigen Alphabet $\Sigma=\{a\}$
\[
 K=\{a^{n^2}\mid n\in\mathbb N_0\}.
\]
Ein Wort gehört also genau dann zu $K$, wenn seine Länge eine Quadratzahl
ist; insbesondere ist $\varepsilon\in K$.
\begin{enumerate}[label=\alph*)]
\item Beweisen Sie mit dem Pumping-Lemma für reguläre Sprachen, dass $K$
nicht regulär ist. Wählen Sie ein geeignetes Wort in Abhängigkeit von der
Pumping-Länge $p$. Argumentieren Sie für \emph{jede} zulässige Zerlegung
und begründen Sie, warum die gepumpte Wortlänge keine Quadratzahl ist.
\punkte{8}
\item Jemand wählt $w=a^{p^2}$, legt die Zerlegung
$x=\varepsilon$, $y=a^p$, $z=a^{p^2-p}$ fest und zeigt, dass
$xy^2z\notin K$ gilt. Warum ist damit allein noch kein vollständiger
Nichtregularitätsbeweis geführt? Geben Sie die Quantorenreihenfolge für
Pumping-Länge, Wort, Zerlegung und Pumpfaktor an, die ein solcher
Widerspruchsbeweis erfüllen muss. \punkte{4}
\end{enumerate}

\section*{Aufgabe 3: Verständnisfragen \hfill 18 Punkte}
\begin{enumerate}[label=\alph*)]
\item \textbf{Operationen auf Sprachen.} Gilt für beliebige Sprachen $A,B$
stets $AB=BA$? Begründen Sie Ihre Antwort durch einen Beweis oder ein
Gegenbeispiel. Entscheiden Sie außerdem, ob die Vereinigung einer regulären
und einer nicht regulären Sprache regulär sein kann. Geben Sie ein
konkretes Beispiel mit Begründung an. \punkte{4}

\item \textbf{Chomsky-Hierarchie.} Beurteilen Sie die folgenden Aussagen
und begründen oder widerlegen Sie jede davon:
\begin{enumerate}[label=\roman*),itemsep=2pt]
\item ,,Eine Sprache, die durch eine kontextfreie Grammatik erzeugt wird,
kann nicht regulär sein.''
\item ,,Jede kontextsensitive Sprache ist auch kontextfrei.''
\item ,,Für jede Sprache vom Typ 0 gibt es eine Turingmaschine, die auf
jeder Eingabe hält und die Zugehörigkeit zur Sprache entscheidet.''
\end{enumerate}
\punkte{6}

\item \textbf{Komplementieren.} Bei einem vollständigen DFA werden
Endzustände und Nichtendzustände vertauscht. Welche Sprache erkennt der
neue Automat? Funktioniert dasselbe Verfahren für beliebige NFA?
Begründen Sie beide Antworten; geben Sie im NFA-Fall gegebenenfalls ein
konkretes Gegenbeispiel an. \punkte{4}

\item \textbf{Zustandsäquivalenz.} Wann sind zwei Zustände eines DFA
äquivalent? Reicht es aus, dass beide keine Endzustände sind?
Begründen Sie Ihre Antwort durch ein Beispiel mit unterscheidender
Fortsetzung. Können zwei äquivalente, vollständige und minimale DFA
unterschiedlich viele Zustände haben? Begründen Sie kurz. \punkte{4}
\end{enumerate}

\ifloesungen
\newpage
\begingroup\small
\setlength{\parskip}{4pt}
\begin{center}{\LARGE\bfseries Musterlösungen}\end{center}
\section*{Lösung 1a: NFA}
Wir wählen $N=(Q,\Sigma,\delta_N,q_0,F_N)$ mit
$Q=\{q_0,\ldots,q_9\}$, $\Sigma=\{a,b\}$ und $F_N=\{q_3,q_6,q_9\}$.
Die drei Ketten lesen $aba$ über $q_1,q_2,q_3$, $bba$ über $q_4,q_5,q_6$
und $bab$ über $q_7,q_8,q_9$.
\begin{center}
\begin{tabular}{c|cc}
 & $a$ & $b$\\\hline
$\to q_0$ & $\{q_0,q_1\}$ & $\{q_0,q_4,q_7\}$\\
$q_1$ & $\varnothing$ & $\{q_2\}$\\
$q_2$ & $\{q_3\}$ & $\varnothing$\\
$*q_3$ & $\varnothing$ & $\varnothing$\\
$q_4$ & $\varnothing$ & $\{q_5\}$\\
$q_5$ & $\{q_6\}$ & $\varnothing$\\
$*q_6$ & $\varnothing$ & $\varnothing$\\
$q_7$ & $\{q_8\}$ & $\varnothing$\\
$q_8$ & $\varnothing$ & $\{q_9\}$\\
$*q_9$ & $\varnothing$ & $\varnothing$
\end{tabular}
\end{center}
In $q_0$ gibt es bei $a$ zwei und bei $b$ drei mögliche Folgezustände.
Der NFA rät dort den Beginn eines Endmusters. Jeder akzeptierende Lauf
liest zuerst ein beliebiges Präfix in $q_0$ und danach genau eines der drei
Muster. Da die Endzustände keine ausgehenden Übergänge haben, muss damit
die Eingabe enden. Umgekehrt lässt sich jedes Wort aus $L$ genau so lesen.

\textit{Bewertung (10):} Zustände, Alphabet, Start- und Endzustände 2;
Übergänge gemäß Bauvorgabe 5; Nichtdeterminismus 1; Korrektheit in beiden
Richtungen 2.

\section*{Lösung 1b: Potenzmengenkonstruktion}
Es gilt $\delta_D(T,c)=\bigcup_{q\in T}\delta_N(q,c)$.
Startmenge ist $\{q_0\}$; eine Menge akzeptiert genau dann, wenn sie
$F_N$ schneidet. Es ergeben sich genau die folgenden neun Zustände:
\begin{center}\small
\begin{tabular}{c|l|cc|l}
 & Zustandsmenge & $a$ & $b$ & erreichbar durch\\\hline
$\to D_0$ & $\{q_{0}\}$ & $D_1$ & $D_2$ & $\varepsilon$\\
$D_1$ & $\{q_{0},q_{1}\}$ & $D_1$ & $D_3$ & $a$\\
$D_2$ & $\{q_{0},q_{4},q_{7}\}$ & $D_4$ & $D_5$ & $b$\\
$D_3$ & $\{q_{0},q_{2},q_{4},q_{7}\}$ & $D_6$ & $D_5$ & $ab$\\
$D_4$ & $\{q_{0},q_{1},q_{8}\}$ & $D_1$ & $D_7$ & $ba$\\
$D_5$ & $\{q_{0},q_{4},q_{5},q_{7}\}$ & $D_8$ & $D_5$ & $bb$\\
$*D_6$ & $\{q_{0},q_{1},q_{3},q_{8}\}$ & $D_1$ & $D_7$ & $aba$\\
$*D_7$ & $\{q_{0},q_{2},q_{4},q_{7},q_{9}\}$ & $D_6$ & $D_5$ & $bab$\\
$*D_8$ & $\{q_{0},q_{1},q_{6},q_{8}\}$ & $D_1$ & $D_7$ & $bba$\\
\end{tabular}
\end{center}
Beispielsweise ist
$\delta_D(\{q_0,q_2,q_4,q_7\},a)
=\{q_0,q_1\}\cup\{q_3\}\cup\varnothing\cup\{q_8\}
=\{q_0,q_1,q_3,q_8\}=D_6$.
Die Tabelle ist unter beiden Übergängen abgeschlossen. Da $q_0$ stets
enthalten ist, wird die leere Menge nicht erreicht.
Der DFA hat Startzustand $D_0$ und Endzustandsmenge $\{D_6,D_7,D_8\}$.

\textit{Bewertung (12):} Zustandsmengen 3; Übergänge 4; Start- und
Endzustände 2; Erreichbarkeitswörter 2; vorgerechnete Vereinigung 1.

\newpage
\section*{Lösung 1c: Minimierung}
Alle neun Zustände sind erreichbar, wie die Wörter in Lösung 1b belegen.
Wir markieren synchron: In Runde $r$ werden nur Markierungen aus den
Runden $0,\ldots,r-1$ verwendet.

\textbf{Anfangsmarkierung ($\times_0$).} Markiert werden die 18 Paare mit
einem Zustand aus $\{D_0,\ldots,D_5\}$ und einem aus $\{D_6,D_7,D_8\}$.
Sie werden durch $\varepsilon$ unterschieden.

\textbf{Runde 1 ($\times_1$).} Bei den nichtakzeptierenden Zuständen bilden
$\{D_0,D_1,D_2\}$, $\{D_3,D_5\}$ und $\{D_4\}$ drei Gruppen:
Aus der ersten Gruppe erreicht kein einzelnes Zeichen einen Endzustand,
aus der zweiten tut dies $a$, aus der dritten $b$. Alle elf Paare zwischen
diesen Gruppen werden markiert. Bei den Endzuständen werden außerdem
$\{D_6,D_7\}$ und $\{D_7,D_8\}$ markiert, etwa mit dem Zeichen $a$.
Es gibt somit 13 neue Markierungen.

\textbf{Runde 2 ($\times_2$).} Die drei Paare innerhalb
$\{D_0,D_1,D_2\}$ werden markiert:
\begin{center}
\begin{tabular}{c|c|c|c}
Paar & Zeichen & bereits markiertes Folgepaar & unterscheidendes Wort\\\hline
$\{D_0,D_1\}$ & $b$ & $\{D_2,D_3\}$ & $ba$\\
$\{D_0,D_2\}$ & $a$ & $\{D_1,D_4\}$ & $ab$\\
$\{D_1,D_2\}$ & $a$ & $\{D_1,D_4\}$ & $ab$
\end{tabular}
\end{center}
\begin{center}
\begin{tabular}{c|cccccccc}
 & $D_0$ & $D_1$ & $D_2$ & $D_3$ & $D_4$ & $D_5$ & $D_6$ & $D_7$\\\hline
$D_1$ & $\times_{2}$ &  &  &  &  &  &  & \\
$D_2$ & $\times_{2}$ & $\times_{2}$ &  &  &  &  &  & \\
$D_3$ & $\times_{1}$ & $\times_{1}$ & $\times_{1}$ &  &  &  &  & \\
$D_4$ & $\times_{1}$ & $\times_{1}$ & $\times_{1}$ & $\times_{1}$ &  &  &  & \\
$D_5$ & $\times_{1}$ & $\times_{1}$ & $\times_{1}$ & $-$ & $\times_{1}$ &  &  & \\
$D_6$ & $\times_{0}$ & $\times_{0}$ & $\times_{0}$ & $\times_{0}$ & $\times_{0}$ & $\times_{0}$ &  & \\
$D_7$ & $\times_{0}$ & $\times_{0}$ & $\times_{0}$ & $\times_{0}$ & $\times_{0}$ & $\times_{0}$ & $\times_{1}$ & \\
$D_8$ & $\times_{0}$ & $\times_{0}$ & $\times_{0}$ & $\times_{0}$ & $\times_{0}$ & $\times_{0}$ & $-$ & $\times_{1}$\\
\end{tabular}
\end{center}
Das Zeichen $-$ bedeutet unmarkiert. Nur die Felder unterhalb der Diagonale
stehen für Paare. Übrig bleiben $D_3\sim D_5$ und $D_6\sim D_8$:
Die jeweiligen Partner haben denselben Akzeptanzstatus und identische
Übergänge. Weitere Markierungen entstehen nicht.

Die Äquivalenzklassen sind
\[
\begin{gathered}
K_0=\{D_0\},\ K_1=\{D_1\},\ K_2=\{D_2\},\ K_3=\{D_3,D_5\},\\
K_4=\{D_4\},\ K_5=\{D_6,D_8\},\ K_6=\{D_7\}.
\end{gathered}
\]
\begin{center}
\begin{tabular}{c|cc@{\qquad}l}
 & $a$ & $b$ & erreichbar durch\\\hline
$\to K_0$ & $K_1$ & $K_2$ & $\varepsilon$\\
$K_1$ & $K_1$ & $K_3$ & $a$\\
$K_2$ & $K_4$ & $K_3$ & $b$\\
$K_3$ & $K_5$ & $K_3$ & $ab$\\
$K_4$ & $K_1$ & $K_6$ & $ba$\\
$*K_5$ & $K_1$ & $K_6$ & $aba$\\
$*K_6$ & $K_5$ & $K_3$ & $bab$
\end{tabular}
\end{center}
Der Quotientenautomat ist vollständig, hat Startzustand $K_0$ und
Endzustandsmenge $\{K_5,K_6\}$. Alle sieben Klassen sind erreichbar.
Zwischen verschiedenen Klassen sind sämtliche Paare markiert, also durch
eine Fortsetzung unterscheidbar. Damit ist der DFA minimal.

\textit{Bewertung (8):} Markierungstabelle mit Runden 4; Klassen 1;
minimaler DFA 2; Erreichbarkeit und Minimalität 1.

\newpage
\section*{Lösung 2: Pumping-Lemma}
\subsection*{a) Nichtregularität}
Angenommen, $K$ sei regulär. Dann gibt es eine Pumping-Länge $p\geq1$,
sodass jedes $w\in K$ mit $|w|\geq p$ eine Zerlegung $w=xyz$ besitzt mit
\[
 |xy|\leq p,\qquad |y|\geq1,\qquad
 \forall i\in\mathbb N_0:\ xy^iz\in K.
\]
Wir wählen
\[
 w=a^{p^2}\in K,\qquad |w|=p^2\geq p.
\]
Sei $w=xyz$ eine \emph{beliebige} Zerlegung mit $|xy|\leq p$ und
$|y|\geq1$. Da das Alphabet nur aus $a$ besteht, gilt $y=a^k$ für ein
$1\leq k\leq p$. Für $i=2$ erhalten wir
\[
 xy^2z=a^{p^2+k}.
\]
Die Länge dieses Wortes liegt strikt zwischen zwei aufeinanderfolgenden
Quadratzahlen:
\[
 p^2<p^2+k\leq p^2+p<p^2+2p+1=(p+1)^2.
\]
Also ist $p^2+k$ keine Quadratzahl und $xy^2z\notin K$.
Dies gilt für jede zulässige Zerlegung und widerspricht dem Pumping-Lemma.
Folglich ist $K$ nicht regulär. \hfill$\square$

Entscheidend ist, dass zwischen den benachbarten Quadratzahlen genug Platz
liegt: Der pumpbare Abschnitt hat höchstens Länge $p$, während die nächste
Quadratzahl den Abstand $2p+1$ besitzt.

\textit{Bewertung (8):} Widerspruchsannahme und Pumping-Länge 1;
geeignetes Wort mit Zugehörigkeit und Länge 2; beliebige Zerlegung und
Schranke für $|y|$ 2; Pumpfaktor und Ausschluss einer Quadratzahl 2; Schluss 1.
Andere vollständige Pumping-Lemma-Beweise sind ebenfalls anzuerkennen.

\subsection*{b) Quantoren und Wahl der Zerlegung}
Die angegebene Zerlegung ist zulässig: $|xy|=p$ und $|y|=p\geq1$.
Auch $p^2+p$ liegt zwischen $p^2$ und $(p+1)^2$.
Trotzdem wurde nur \emph{eine} Zerlegung widerlegt. Für eine reguläre Sprache
muss das Pumping-Lemma lediglich die Existenz \emph{einer geeigneten}
Zerlegung garantieren. Es könnte also eine andere Zerlegung geben, für die
alle Pumpfaktoren funktionieren. Genau das muss der Beweis ausschließen.

Die nötige Quantorenreihenfolge für den Widerspruch zur Pumping-Eigenschaft
lautet:
\[
 \forall p\geq1\ \exists w\in K,\ |w|\geq p\
 \forall\text{ zulässigen Zerlegungen }w=xyz\
 \exists i\in\mathbb N_0:\ xy^iz\notin K.
\]
Das Wort darf von $p$ abhängen; der Pumpfaktor darf von der Zerlegung
abhängen. In a) funktioniert sogar einheitlich $i=2$.
Die Zerlegung selbst darf man nicht auf einen günstigen Spezialfall
beschränken.

\textit{Bewertung (4):} Unterschied zwischen einer und allen zulässigen
Zerlegungen 2; richtige Quantorenreihenfolge einschließlich der
Zulässigkeitsbedingungen (wie in a) definiert) 2.

\newpage
\section*{Lösung 3: Verständnisfragen}
\subsection*{a) Operationen auf Sprachen (4 Punkte)}
\textbf{Konkatenation ist nicht kommutativ.} Für $A=\{a\}$ und $B=\{b\}$
gilt $AB=\{ab\}$, aber $BA=\{ba\}$; diese Sprachen sind verschieden.
\textit{(2 Punkte)}

\textbf{Die Vereinigung kann regulär sein.} Über $\Sigma=\{a,b\}$ wählen wir
$R=\Sigma^*$ und $N=\{a^nb^n\mid n\in\mathbb N_0\}$.
$R$ ist regulär, $N$ ist bekanntlich nicht regulär, aber
$R\cup N=\Sigma^*$ ist regulär. Nichtregularität eines Operanden erzwingt
also nicht die Nichtregularität der Vereinigung. \textit{(2 Punkte)}

\subsection*{b) Chomsky-Hierarchie (6 Punkte)}
\textbf{i) Falsch.} Jede reguläre Sprache ist auch kontextfrei.
Beispielsweise ist $S\to aS\mid\varepsilon$ eine kontextfreie Grammatik
für die reguläre Sprache $a^*$. Dass eine Grammatik kontextfrei ist,
schließt eine reguläre Beschreibung ihrer Sprache nicht aus.
\textit{(2 Punkte)}

\textbf{ii) Falsch.} Die Inklusion verläuft umgekehrt: Jede kontextfreie
Sprache ist kontextsensitiv (mit der üblichen Sonderkonvention für
$\varepsilon$), aber nicht jede kontextsensitive Sprache ist kontextfrei.
Ein Gegenbeispiel ist $\{a^nb^nc^n\mid n\geq1\}$.
\textit{(2 Punkte)}

\textbf{iii) Falsch.} Typ-0-Sprachen sind rekursiv aufzählbar bzw.\
semi-entscheidbar. Eine erkennende Turingmaschine akzeptiert jedes Wort der
Sprache, darf bei Wörtern außerhalb aber unendlich weiterlaufen.
Nicht jede solche Sprache ist entscheidbar; ein Beispiel ist die Sprache
des Halteproblems. \textit{(2 Punkte)}

\subsection*{c) Komplementieren (4 Punkte)}
\textbf{DFA:} Der neue Automat erkennt $\Sigma^*\setminus L(M)$.
Für jedes Wort gibt es genau einen vollständigen Lauf. Dessen letzter
Zustand ist nach dem Vertauschen genau dann akzeptierend, wenn er vorher
nicht akzeptierend war. Vollständigkeit sichert, dass der Lauf für jedes
Wort existiert. \textit{(2 Punkte)}

\textbf{NFA:} Das Verfahren funktioniert im Allgemeinen nicht. Ein NFA
akzeptiert, sobald \emph{mindestens ein} Lauf akzeptiert. Zum Verwerfen
müssen hingegen \emph{alle} Läufe nichtakzeptierend sein.
Als Gegenbeispiel wähle man $Q=\{s,f\}$, $\Sigma=\{a\}$, Startzustand $s$,
$F=\{f\}$ sowie $\delta(s,a)=\{s,f\}$ und $\delta(f,a)=\{f\}$.
Das Wort $a$ wird über den Lauf nach $f$ akzeptiert. Nach dem Vertauschen
ist $F'=\{s\}$; dann wird $a$ über den Lauf nach $s$ weiterhin akzeptiert.
Das kann nicht das Komplement sein. \textit{(2 Punkte)}

\subsection*{d) Zustandsäquivalenz (4 Punkte)}
Zwei Zustände $q,r$ sind äquivalent, wenn für jede Fortsetzung
$v\in\Sigma^*$ gilt
\[
 \delta^*(q,v)\in F\quad\Longleftrightarrow\quad\delta^*(r,v)\in F.
\]
Hier bezeichnet $\delta^*$ die auf Wörter erweiterte Übergangsfunktion.
\textit{(1 Punkt)}

Gleicher Akzeptanzstatus allein genügt nicht: Im DFA über $\{a\}$ mit
Zuständen $q,r,f$, Startzustand $q$, $F=\{f\}$ und Übergängen
$q\xrightarrow{a}r$, $r\xrightarrow{a}f$, $f\xrightarrow{a}f$
sind $q$ und $r$ beide nichtakzeptierend. Die Fortsetzung $a$ führt von
$q$ in einen Nichtendzustand, von $r$ dagegen in einen Endzustand.
\textit{(2 Punkte)}

Zwei äquivalente, vollständige und minimale DFA über demselben Alphabet
haben gleich viele Zustände: Beide haben die kleinstmögliche Zustandszahl
für dieselbe Sprache. Sie sind sogar bis auf Umbenennung der Zustände
isomorph. \textit{(1 Punkt)}
\endgroup
\fi
\end{document}
